%% LyX 2.5.2 created this file.  For more info, see https://www.lyx.org/.
%% Do not edit unless you really know what you are doing.
\documentclass[a4paper,ngerman,year=2024, month=April]{article}
\PassOptionsToPackage{obeyFinal}{todonotes}
\usepackage{fontspec}
\setcounter{secnumdepth}{3}
\setcounter{tocdepth}{3}
\usepackage{enotez}
\usepackage{todonotes}
\usepackage[authoryear]{natbib}

\makeatletter

%%%%%%%%%%%%%%%%%%%%%%%%%%%%%% LyX specific LaTeX commands.
% Backwards compatibility for LuaTeX < 0.90
\@ifundefined{pageheight}{\let\pageheight\pdfpageheight}{}
\@ifundefined{pagewidth}{\let\pagewidth\pdfpagewidth}{}
\pageheight\paperheight
\pagewidth\paperwidth

\providecommand\textquotedblplain{%
  \bgroup\addfontfeatures{RawFeature=-tlig}\char34\egroup}

%%%%%%%%%%%%%%%%%%%%%%%%%%%%%% Textclass specific LaTeX commands.
\usepackage{maa-monthly-todo}
\numberwithin{equation}{section}
\numberwithin{figure}{section}

%%%%%%%%%%%%%%%%%%%%%%%%%%%%%% User specified LaTeX commands.
\RequirePackage{scrextend}

\usepackage[
    colorlinks=true,    % Aktiviert farbigen Text (schaltet die Rahmen ab)
    linkcolor=blue,     % Farbe für interne Links (z. B. Inhaltsverzeichnis, Abbildungen)
    citecolor=blue,     % Farbe für Literaturzitate im Text
    urlcolor=blue,      % Farbe für URLs/Webseiten im Literaturverzeichnis
    breaklinks=true     % Erlaubt Links, am Zeilenende umzubrechen
]{hyperref}
\usepackage{url}

\usepackage{esint}
\usepackage{lipsum}

\usepackage{pgfplots}
\pgfplotsset{compat=1.18}

\usepackage{siunitx}
\sisetup{locale = DE} % Uses comma as decimal separator
%\sisetup{inter-unit-product = \cdot} % Mittelpunkt zwischen Einheiten
% Das Leerzeichen komplett entfernen
%\sisetup{inter-unit-product = {}, separate-uncertainty = true} 
%\sisetup{inter-unit-product = \,} % Standard (schmales Leerzeichen)

% \DeclareSIUnit\kwh{kWh} % Eigendefinition: Die Energie beträgt \qty{0,6}{\kwh}.

\usepackage{booktabs}

\usepackage{wasysym}

%% List of urls
\usepackage{tocloft}
\usepackage{xurl}

% Define a custom list structure for links
\newcommand{\listlinksname}{List of Links}
\newlistof{links}{lnks}{\listlinksname}

% Intercept the standard \href command used by LyX's Hyperlink inset
\let\oldhref\href
\renewcommand{\href}[2]{\oldhref{#1}{#2}\addcontentsline{lnks}{figure}{#2: \url{#1}}}


%----------- Riviera ist schuld -----------

\usepackage{luacode}
\begin{luacode*}
local function uebersetze_einheiten(text)
    -- Ersetzt "tb" nach einer Zahl durch "EL"
    text = string.gsub(text, "(%d+)%s*tb", "%1 EL")
    -- Ersetzt "tsp" nach einer Zahl durch "TL"
    text = string.gsub(text, "(%d+)%s*tsp", "%1 TL")
    return text
end

-- Aktiviert die Ersetzung für das gesamte Dokument
luatexbase.add_to_callback("process_input_buffer", uebersetze_einheiten, "Uebersetzer")
\end{luacode*}

\makeatother

\providecommand\theoremname{Theorem}
\theoremstyle{plain}
\newtheorem{thm}{\protect\theoremname}
\providecommand\claimname{Behauptung}
\theoremstyle{remark}
\newtheorem{claim}[thm]{\protect\claimname}
\usepackage{polyglossia}
\setdefaultlanguage[variant=german,spelling=new,babelshorthands=true]{german}
\begin{document}
\makeatletter
% Holt die globalen Optionen der article-Klasse und prüft, ob "twoside" drinsteht
\@ifclasswith{article}{twoside}{%
  % Dieser Code wird ausgeführt, wenn twoside in den Dokumentenklassen-Optionen aktiv ist
  \newcommand{\twosidedmode}{1} 
}{%
  % Dieser Code wird ausgeführt, wenn twoside NICHT aktiv ist
  \newcommand{\twosidedmode}{0} 
}
\makeatother


\renewcommand\textHeadAssoc{Käthe Kolwitz}
\renewcommand\textHeadJournal{Jedzia Monthly}
\renewcommand\textFootAssoc{Mälscher Brieh}

% Setzt das Jahr (z.B. 2024)
\setcounter{annual}{2026}

% Setzt den Monat (1=January, 2=February, 3=March, 4=April, ...)
\setcounter{issue}{10} 

\makeatletter

% Forces a larger base font size (e.g., 14pt) and scales everything up
\changefontsizes[16pt]{14pt} % 12pt document size, 14pt line spacing
%\changefontsizes[12pt]{11pt} % 12pt document size, 14pt line spacing

% Wenn zweiseitig mit angepasstem Layout und Randnoten
\ifnum\twosidedmode=1
 % Wenn 1: Aktivierte das zweiseitige Layout und passe die Ränder an
  %\geometry{twoside, outer=4cm, inner=2cm, marginparwidth=3cm, marginparsep=3mm}

%\hsize26pc\textwidth26pc
%\hsize152mm
%\setlength\textwidth{139mm}      % 210mm - 2*26mm = 158mm Textbreite
\setlength\oddsidemargin  {2mm}
\setlength\evensidemargin {18mm}
% Seitenrand
%\setlength{\marginparsep}{3mm} % Verringert den Abstand zum Text auf 3 Millimeter
\setlength{\marginparsep}{5mm} % Verringert den Abstand zum Text auf 3 Millimeter
% Todo Notes
\renewcommand{\@todonotes@textwidth}{2cm}
% 1. Die physische Breite für alle Randnotizen (\marginpar) auf 4cm setzen
\setlength{\marginparwidth}{2cm}



%\else

%\setlength{\marginparsep}{3mm} % Verringert den Abstand zum Text auf 3 Millimeter
\fi




% nach section neue Zeile
\renewcommand\section{\@startsection{section}{1}{\z@}%
  {-3.5ex \@plus -1ex \@minus -.2ex}% % Abstand vor der Überschrift
  {2.3ex \@plus .2ex}%                 % Wert > 0 erzwingt den Zeilenumbruch danach
  {\normalfont\Large\bfseries}}

% nach subsection neue Zeile
\renewcommand\subsection{\@startsection{subsection}{2}{\z@}%
  {-3.25ex\@plus -1ex \@minus -.2ex}%
  {1.5ex \@plus .2ex}% % Dieser Wert > 0 erzwingt den Zeilenumbruch
  {\normalfont\large\bfseries}}

%\renewcommand\subsubsection{\@startsection{subsubsection}{3}{\z@}%
%  {-12pt \@plus -2pt \@minus -2pt}%
%  {1.5ex \@plus .2ex}% % Dieser Wert > 0 erzwingt den Zeilenumbruch
%  {\normalfont\small\itshape}}
\renewcommand\subsubsection{\@startsection{subsubsection}{3}{\z@}%
                                     {-12pt \@plus -2pt \@minus -2pt}%
                                     {-1.5em}%
%                                     {\normalfont\normalsize\itshape}}
  {\normalfont\normalsize\bfseries\itshape}}

\renewcommand\paragraph{\@startsection{paragraph}{3}{\z@}%
  {-12pt \@plus -2pt \@minus -2pt}%
  {-.5em}%
  {\normalfont\normalsize\itshape}}

\renewcommand\subparagraph{\@startsection{subparagraph}{3}{5mm}%
  %{-12pt \@plus -2pt \@minus -2pt}%
  %{-3pt \@plus -1pt \@minus -1pt}% <-- Hier wird der Abstand nach oben verringert
  {-0pt \@plus  0pt \@minus  0pt}% ganz abschalten
  {-.5em}%
  {\normalfont\normalsize\itshape}}
%\renewcommand\subparagraph{\@startsection{subparagraph}{5}{\z@}%
%  {3.25ex \@plus1ex \@minus.2ex}%
%  {-1em}%
%  {\normalfont\small\itshape}}
%\renewcommand\subparagraph{\@startsection{subparagraph}{5}{5mm}%
%  {-12pt \@plus -3pt \@minus -2pt}% <-- Abstand VOR der Überschrift
%  {-0.5em}%                          <-- Abstand NACH der Überschrift (Inline)
%  {\normalfont\small\itshape}}

\makeatother


% Endnotes Einstellungen
\makeatletter
% --- VARIABLEN-DEFINITION ---
\newcommand{\jinfo}{\textcircled{\small i}}
% Hier das gewünschte Zeichen eintragen (z. B. e, n, oder x)
%\newcommand{\enoteprefix}{A}
%\newcommand{\enoteprefix}{\bell}
\newcommand{\enoteprefix}{\jinfo}
% ----------------------------


% Ein eigenes Makro definieren, das das 'e' mit in das Argument zieht
% 1. Formatierung für den Fließtext (hochgestellt & verlinkt)
% Verwendet die Variable \enoteprefix vor dem Argument #1
\newcommand{\mycustommark}[1]{\textsuperscript{\enoteprefix#1}}

\setenotez{
  % Weist enotez an, unser Makro anstelle von \textsuperscript zu nutzen
  %mark-cs = \mycustommark, 
  mark-format = Anm:\ ,
  %counter-format = Alph
}

% 2. Formatierung für das Verzeichnis (Explizite Definition mit 1 Parameter)
%\renewcommand{\enmark}[1]{\enoteprefix#1\quad}
\makeatother


% Rule with Text
\makeatletter
\newcommand{\linewithtext}[1]{%
  \par\noindent\raisebox{.5ex}{\makebox[\linewidth]{\hrulefill\hspace{1ex}\raisebox{-.5ex}{#1}\hspace{1ex}\hrulefill}}\par
}
\makeatother

\title{Muscheln}
\author{von Stefanie Heinzmann}
\date{3. Oktober 2026}
\markright{Schalentiere}
\maketitle
\begin{abstract}
Test.\todo{Idee von Riviera. {*}Überkompliziert{*} aber gut}
\end{abstract}
\begin{claim}
MMMMM-{}-{}-{}-{}- Recipe via Meal-Master (tm) v8.06, „Netzmarder“
wird hier nicht als Wort verändert!

Sondern alle „tb“'s in Essloeffel transferiert! (EL).
\end{claim}


\section{Muscheln in Wuerzsud}

Categories: Weichtier, Salzwasser, Muschel

Yield: 4 Servings

3 kg Muscheln

250 g Zwiebeln

3 Knoblauchzehen

2 Fleischtomaten; enthaeutet

2 bn Petersilie

40 g Butter

Cayennepfeffer

1/4 l Weisswein

1 Zitrone; Saft

Muscheln in reichlich kaltes Wasser geben, einige Stunden darin liegen

lassen; das Wasser ab und zu erneuern, die Muscheln anschliessend

gruendlich buersten. Bartbueschel entfernen, so lange spuelen bis
das

Wasser vollkommen klar bleibt.

Muscheln die sich beim Waessern und Buersten oeffnen, sind

ungeniessbar, nur geschlossene Muscheln verwenden.

Zwiebeln in Scheiben schneiden, Knoblauchzehen wuerfeln, enthaeutete

Fleischtomaten grob zerkleinern (die Stengelansaetze herausschneiden),

Petersilie fein hacken.

Butter in einem grossen breiten Topf zerlassen; Zwiebel und Knoblauch

darin anduensten, Tomatenstuecke und Petersilie unterruehren.

Cayennepfeffer, Weisswein, Zitronensaft hinzugiessen, erhitzen, die

Muscheln hinzufuegen, zugedeckt bei starker Hitze zum kochen bringen.

10-12 Minuten kochen lassen, dabei den Topf kraeftig ruetteln, damit

die Muscheln gleichmaessig garen.

Erfasst und gepostet von Klaus Wehnert@2:2457/145.10, 29.04.94,

textlich etwas angepasst

\section{Miesmuschelnsuppe 'Chez Benoit'}

Categories: Suppen

Yield: 4 Servings

1 kg Frische Miesmuscheln

1 1/4 dl Weisswein, trocken (z.B.

-{}- Muscadet)

1 1/4 dl Rahm

2 tb Butter

1 lg Zwiebel, in duennen Ringe

2 sm Moehren, in duennen Scheiben

2 Stuecke Lauch (nur den

-{}- weissen Teil nehmen),

-{}- in duennen Ringen

100 g Speck, fein geschnitten

100 g Geraeuchertes Schinken,

-{}- fein geschnitten

5 dl Fisch- oder Gefluegelfond

100 g Blanchierte Bohnen, schraeg

-{}- geschnitten, Stuecke von

-{}- ca. 15 mm Laenge (1)

Blanchieren: Die Bohnen 4 Minuten in reichlich kochendem Salzwasser

erwaermen, anschliessend mit kaltem Wasser so lange abspuelen, bis
die

Bohnen kalt sind.

Die Muscheln gut buersten und waschen, entbarten (nicht im

voraus entbarten, sonst sterben die Muscheln ab!!).

Weisswein und Rahm in einer Pfanne tun, Muscheln zugeben und

aufkochen (die Muscheln sollten knapp gedeckt sein: falls nicht der

Fall, Wein und Rahm zu gleichen Mengen zugeben). Dabei den Topf

immer gut ruetteln; nur so lange kochen, bis die Muscheln auf

sind (ca. 3 bis 4 Minuten; nicht zu lange kochen, sonst werden die

Muscheln zaeh).

Muscheln abtropfen lassen (Kochfluessigkeit auffangen), die

noch geschlossenen Muscheln WEGWERFEN. Die Kochfluessigkeit sehr gut

filtrieren und beiseite tun. Muschelfleisch aus den Schalen nehmen,

beiseite tun.

Butter in einer Pfanne aufschmelzen, Zwiebel, Moehren, Lauch,

Speck und Schinken darin duensten, ca. 5 Minuten (die Zwiebeln

duerfen nicht braun werden !). Fischfond zugeben, koecheln lassen,

bis die Gemuese sehr weich sind (ca. 45 Minuten).

Muscheln-Kochfluessigkeit zugeben, aufkochen.

5 Minuten vor dem Servieren, Muschelfleisch und Bohnen zugeben,

koecheln lassen, bis wieder warm und mit Salz und Pfeffer

abschmecken. Sofort servieren.

Aus: Patricia Wells, Bistro cooking, Arrow London 1992, ISBN

0-09-992340-8

09.12.1993

\section{Toskanische Fischsuppe (Cacciucco alla toscana)}

Categories: Suppen, Fisch

Yield: 99 Servings

1 1/2 kg Seefische, wie Makrelen,

-{}- Seehecht, Goldbrassen,

-{}- Merlan, Schellfisch,

-{}- Barsch, Sardinen

10 Krebse oder Garnelen

1 kg Muscheln

3 Knoblauchzehen (oder mehr)

1 Zwiebel

1 sm Peperoni

500 g Dosen-Tomaten (keine

-{}- Hollandtomaten)

1/4 l Weisswein z.B. Frascati

1 ct Fleischsuppe oder die

-{}- entsprechende Menge Fond

1/2 ts Paprika

1/4 ts Thymian, am besten frisch

-{}- (ca. ein Zweig)

Olivenoel

4 sl (-8) Weissbrot oder Baguette

Die Fischsorten - moeglichst 4-6 verschiedene - kuechenfertig

vorbereiten und in Stuecke schneiden. Die Muscheln und Krebse

unter fliessend Wasser gruendlich abbuersten. In einem grossen Topf

ca. 8 EL Olivenoel erhitzen. Eine Knoblauchzehe, die Zwiebel, die

Petersilie und die entkernte Peperoni feinhacken und in dem Oel

anduensten. Wenn frische Tomaten verwendet werden diese

ueberbrueht und enthaeutet. Die Tomaten puerieren -

moeglichst ohne Kerne durch ein Sieb passieren - und in den Topf

geben. Die Fleischbruehe - man kann auch Fischfond nehmen -

und den Wein dazugiessen. Alles ca. 15 min kochen lassen.

Nun die Fischstuecke dazugeben, die Muscheln und Krebse zuerst.

Jetzt das ganze nochmal ca. 15 bis 20 min kochen lassen und dann mit

Peffer, Salz, Thymian und Paprika abschmecken. Die

Weissbrotscheiben mit den durchgeschnittenen Knoblauchzehen

einreiben, in heissem Olivenoel goldbraun backen oder im Backofen

ohne Fett roesten. Die Scheiben in tiefe Suppenteller oder eine

grosse Terrine legen und darueber die Fischsuppe geben.

Beilagen : Fruchtiger Weisswein

07.01.1994

\section{Sommersalat mit Meeresfruechten}

Categories: Salate, Fisch

Yield: 4 Servings

1 Eisbergsalat

4 Feste Tomaten

4 Hartegekochte Eier

1 Zwiebel

1 Gruene Paprika

1 cn Thunfisch in Oel

1 cn Miesmuscheln

75 g Krabbenfleisch

50 g Doppelrahm-Frischkaese

1 Eigelb

4 tb Heisses Wasser

8 tb Oel

2 tb Weinessig

1 tb Senf

1 Knoblauchzehe

Pfeffer

Salz

Eisbergsalat putzen und in ca. 4 cm grosse Stuecke schneiden, dann

eine flache Salatschuessel damit auslegen. Tomaten in Spalten

schneiden und die Eier achteln. Zwiebeln hacken, Paprika in Streifen

schneiden, alles in das Salatbett geben und den zerbroeckelten

Thunfisch, abgetropfte Muscheln und Krabben darueber verteilen.

Den Frischkaese mit einer Gabel fein zerdruecken. Alle

weiteren Zutaten nach und nach hinzufuegen und zu einer cremigen

Sosse verruehren.

Pro Portion: ca. 2390 kj/570 kcal

Beilage : Baguette

07.01.1994

\section{Gefuellte Gebackene Muscheln}

Categories: Vorspeisen, Fisch

Yield: 8 Servings

20 Frische Zehen Knoblauch,

-{}- oder mehr (!)

3 cn Muscheln a' 180 g,

-{}- abgegossen und gehackt

1 1/2 c Frische Muscheln, gekocht

-{}- und gehackt

150 g Weiche Butter

1 tb Oregano

2 tb Spinat, gekocht und

-{}- ausgedrueckt

4 tb Sherry

150 g Semmelbroesel

1 bn Petersilie, fein gehackt

2 tb Zitronensaft

2 tb Pinienkerne oder gehackte

-{}- Walnuesse

Muskat

Cayenne-Pfeffer

Salz

Pfeffer aus der Muehle

Den Knoblauch fein hacken, durch die Presse geben und in einer

Schuessel mit allen Zutaten bis auf die Pinienkerne und den

Cayenne-Pfeffer mischen. Die Masse in eine grosse, ofenfeste Form

oder portioniert in kleine Formen oder Muschelschalen fuellen, mit

den Pinienkernen (oder Nuessen) sowie Cayenne-Pfeffer bestreuen

und bei 180 Grad im Backofen ueberbacken, bis alles eine

goldbraune Farbe annimmt.

Man kann die Muscheln mit Zitronenschnitzeln und Dill dekoriert

heiss (!) servieren.

08.01.1994

\section{\textcolor{red}{Allgemeines zu Miesmuscheln}}

Categories: Weichtiere

Yield: 1 Anleitung

Muscheln muessen frisch sein, das heisst noch leben, wenn sie in

den Topf kommen. Zoologisch handelt es sich bei Miesmuscheln um

Weichtiere. Aus Italien sind die Muscheln als Cozze oder Vongole

bekannt, in Frankreich heissen sie Moules, in Deutschland und in

der Schweiz (Erntegebiet ist dort die Nordseekueste von Daenemark

bis Holland) sind sie als Miesmuscheln auf dem Markt.

Im Gegensatz zu ihren \textquotedbl Luxusschwestern\textquotedbl{}
- den Austern - werden

die Muscheln selten roh gegessen.

Verdorbene Muscheln sind gluecklicherweise leicht zu erkennen:

Muscheln, die bereits offen sind, sowie Tiere mit beschaedigten

Schalen wegwerfen. Ebenfalls ungeniessbar sind Muscheln, die sich

in der Pfanne nicht von selbst oeffnen. In diesen Faellen sind die

Tiere wahrscheinlich bereits auf dem Transport gestorben: gesunde

Muscheln bleiben mit der Fluessigkeit, die sie im Inneren

speichern koennen, so lange am Leben, bis sie verzehrt werden.

Miesmuscheln werden haeufig kuechenfertig gereinigt angeboten:

die Schalen sind entsandet und der sogenannte Byssusfaden, ein

schnuratiges Gebilde, mit dem sich die Muschel am Untergrund

festhaelt, ist bereits weggeschnitten. Sie sollten dann auf dem

Laden nicht viel aelter als 24 Stunden sein. Ist der Faden noch

dran, mit einem Messer abschneiden (\textquotedbl entbaerten\textquotedbl )
und unter

fliessendem kaltem Wasser gut reinigen. Und dann sofort: ab in die

Pfanne !

12.01.1994

\section{Moules a la mariniere (Grundrezept)}

Categories: Weichtiere

Yield: 4 Servings

2 kg Miesmuscheln

2 tb Olivenoel

1 lg Zwiebel, kleingeschnitten

1 Lorbeerblatt

2 1/2 dl Weisswein

Rosmarin

Thymian

Salz

Pfeffer

MMMMM-{}-{}-{}-{}-{}-{}-{}-{}-{}-{}-{}-{}-{}-{}-{}-{}-{}-{}-{}-{}-{}-{}-{}-{}-SERVIEREN-{}-{}-{}-{}-{}-{}-{}-{}-{}-{}-{}-{}-{}-{}-{}-{}-{}-{}-{}-{}-{}-{}-{}-{}-{}-{}-{}-{}-{}-

50 g Butter

Zitronenspalten

Knoblauchbrot

Miesmuscheln vorbereiten: siehe Anleitung.

Olivenoel in einer hohen Pfanne erhitzen, Zwiebel,

Lorbeerblatt und Kraeuter beigeben, sanft durchdaempfen.

Miesmuscheln zugeben und mit dem Weisswein abloeschen. Mit Salz

und Pfeffer leicht wuerzen und unter sorgfaeltigem Wenden mit

einer Holzkelle oder unter sanftem Schuetteln und Ruehren etwa 8

Minuten auf kleiner Flamme zugedeckt kochen lassen, bis sich die

Muscheln geoeffnet haben (die Muscheln, die sich nicht

oeffnen, wegwerfen).

Die Muscheln in ein Sieb schuetteln, den Sud fuer allfaellige

Weiter- zubereitung auffangen.

Als Moules marinieres servieren:

Die Muscheln in grossen, vorgewaermten Tellern anrichten,

Kraeuter entfernen. Die kalte Butter in kleinen Stueckchen mit dem

Schwingbesen im Fonds ruehren und die Sauce ueber die Muscheln

anrichten. Mit Zitronen- spalten garnieren und mit knusprigem

Knoblauchbrot servieren.

12.01.1994

\section{Muschelrisotto mit Fenchel}

Categories: Weichtiere, Gemuese

Yield: 4 Servings

2 kg Moules a la mariniere, siehe

-{}- Grundrezept

750 g Fenchel, geputzt, klein

-{}- gewuerfelt

-{}- Etwas Fenchelkraut hacken

2 tb Olivenoel

250 g Reis (Arborio, Vialone)

5 dl Bouillon (MENGE ANPASSEN)

50 g Butter

50 g Parmesan, gerieben

Weisser Pfeffer

Die Miesmuscheln nach dem Grundrezept, aber ohne Butter

zubereiten. Die Muscheln in ein Sieb schuetten, den Sud auffangen.

Olivenoel in einer Pfanne erhitzen, Fenchelwuerfel und Reis darin

glasig duensten, dabei ruehren. Muschelsud durch ein Haarsieb

giessen und den Reis damit abloeschen. Den Risotto waehrend ca.

18 Minuten immer wieder sanft ruehren.

Inzwischen die Muscheln, bis auf 12 Stueck, aus den Schalen loesen.

Das Muschelfleisch vor dem Servieren unter der heissen

Risotto mischen; Parmesan, die Butter und das gehackte

Fenchelkraut sanft darunterruehren und auf einer heissen Platte

anrichten. Beiseite gelegte Muscheln in der Schale darauf legen und

sofort servieren. 

12.01.1994

\todo[inline, color=yellow]{Das Ende ...}

\rule[0.5ex]{1\columnwidth}{2pt}

\notizenxl

\listoftodos{}

\printendnotes

\clearpage{}

\hspace{2cm}\linewithtext{Ende des eigentlichen Textes}

\subsubsection*{Notizen:}

Test Test Test, 
\end{document}
